\chapter{Conclusion and Outlook}
\label{ch:conclusion}

\section{Summary}
\label{sec:summary}

\TODO{Answer RQ1 directly in a short paragraph, then summarise the contributions:
an asynchronous, content-addressed ingestion pipeline with layered idempotency;
a retrieval-grounded tutor with four distinct groundedness mechanisms; and a
systematic audit method that separates implemented from delivered behaviour.}

\section{Contributions}
\label{sec:contributions}

\TODO{State them as claims a reader can check, not as activities. "Implemented a
pipeline" is an activity. "Demonstrated that six independent idempotency
mechanisms across four layers make repeated ingestion of identical input free,
including under concurrent workers" is a claim.}

\section{Limitations}
\label{sec:limitations}

\TODO{Consolidate from \cref{sec:threats} and be concrete. The ones that matter:
no learning-outcome study; prompt-level injection defences are mitigations
rather than guarantees; the human review step cannot gate publication; and
several implemented features are undeployed, so their behaviour at real usage is
unobserved.}

\section{Future Work}
\label{sec:future-work}

\TODO{Prefer items this thesis has already shown to be necessary over generic
wishes. Strong candidates, each traceable to a finding:
(1) surface citations in the student interface, closing the gap in
\cref{sec:citation-gap};
(2) introduce a draft/published state so review becomes a real gate
(\cref{sec:review});
(3) decouple the provider chain into genuinely independent failure domains and
classify authentication and payment failures (\cref{sec:provider-strategy});
(4) run the empirical study designed in \cref{sec:unanswered};
(5) generate architecture documentation from code so the drift measured in
\cref{sec:reality-vs-documented} cannot recur.}

\section{Reflection}
\label{sec:reflection}

\TODO{Optional, and only if your institute welcomes it --- some do, some regard
it as padding. If included, the honest reflection worth writing is what the
audit taught about the gap between what a developer believes a system does and
what it does, and why documentation is not evidence.}
